\subsection{Finite arithmetic descent for generalized theories}

The following argument uses the explicit uniform eventual conditions of
scans~117 and~119. We keep both gradings nonnegative and preserve them
throughout. In particular, the construction also preserves the additional
support condition $q\le w$ if that condition is imposed: our paths are
simplicial cotensors, not differential forms in a new weight-zero
variable.

\begin{lemma}[C16: \status{Proved}; paths and marked homotopies]\label{sa08:C16}
Let $k$ be any commutative ring and let $A$ be an admissible generalized
theory. Let $I_1=\Delta[1]$ and let $I_2$ be two consecutive edges, with
vertices $0,1,2$ and edges $01,12$. Define the simplicial cotensor by
\[
 (A^{I_j})_n^{w,q}
 =\Hom_{s\operatorname{Mod}_k}
       (k[\Delta[n]\times I_j],A^{w,q}).
\]
All algebraic structures are defined pointwise. Then $A^{I_j}$ is
admissible, and its endpoint maps and constant-path section preserve
all structures. If $A$ has flat terms, so does $A^{I_j}$.

The relative complex of paths vanishing at the initial endpoint is
vertically acyclic in every fixed cochain range in all sufficiently
large weights, uniformly in $n$. The relative complex vanishing at
both endpoints has, in the same eventual sense, only coefficient
cohomology $S^w$ in degree one.

Suppose $G_N$ is an elementary source of
Statement~\ref{sa08:C09} over a coefficient ring where that source is
defined, and $A^{N,0}$ is simplicially acyclic. Then the space of
marked maps defining $G_N\to A$ is nonempty and contractible. More
explicitly, any compatible boundary family of such maps extends
over its simplex; compatible squares of marked homotopies also
extend over $\Delta[1]\times\Delta[1]$.
\end{lemma}
\begin{proof}
Cotensoring preserves the two gradings and all pointwise operations,
including divided powers. It also preserves augmentations, since
$k^{I_j}=k$. If a simplicial module $M$ is acyclic, restriction along
any inclusion of finite simplicial sets is surjective: the quotient
of their free simplicial modules has degreewise projective normalized
terms, so the Ext computation of Statement~\ref{sa08:C02} applies.
The same normalized Hom calculation shows that $M^{I_j}$ is acyclic.
This proves eventual horizontal acyclicity of the cotensors.

We spell out the uniform comparison used vertically. Fix a cochain
range $0\le q\le d$. Choose a weight threshold beyond the finitely
many resolution and flasqueness bounds through degree $d+2$.
In that range, $S^w\to A^{w,*}$ is a resolution on the simplicial
topos, and its relevant terms are acyclic on every slice by
Statement~\ref{sa08:C02}. The hypercohomology comparison therefore
identifies sections through degree $d$ with coefficient cohomology
on every simplicial set. Indeed, the cone of the augmentation has
cohomology only above the chosen range, while the positive derived
sections of each relevant term vanish. Both are first-quadrant
calculations in the fixed range. Thus the threshold does not depend
on the simplicial set.

Apply this to $\Delta[n]\times I_j$ and to its endpoint subobjects.
The product is contractible. The pair with one endpoint is acyclic,
and the pair with both endpoints has relative coefficient cohomology
$k$ in degree one only. In the chosen weights restriction on each
relevant cochain term is surjective by the first paragraph. The
relative section sequences and their long exact cohomology sequences
give all the asserted vertical comparisons, uniformly in $n$.

For flatness, the finite integral normalized chain complex of
$\Delta[n]\times I_j$ is a copy of $\Z[0]$ plus finite free identity
disks. To justify the splitting, its reduced complex is exact and
finite free; each cycle subgroup is free, and the short exact cycle
sequences split. Tensor the splitting with $k$ and apply chain Hom
into $N(A^{w,q})$. The resulting module of chain maps is a finite
direct sum of normalized terms of $A^{w,q}$, hence flat. For the pair
with the initial endpoint, only identity disks occur. For the pair
with both endpoints, the corresponding splitting is $\Z[1]$ plus
identity disks; its Hom uses $Z_1(N(A^{w,q}))$ and normalized terms.
We use this latter description below only when those cycles are
known to be flat.

For the marked-map assertion write $\Phi_n=k^{n+1}$, with marked
constant vector $1$. Nonemptiness is the normalized disk lifting
in Statement~\ref{sa08:C09}. For a boundary inclusion $L\subset K$,
a prescribed marked family is a linear map on
\[
 (\Phi\otimes k[L])+(k\cdot1\otimes k[K])
 \ \subset\ \Phi\otimes k[K].
\]
Its quotient is $(\Phi/k)\otimes k[K/L]$, which is degreewise free.
Statement~\ref{sa08:C02} makes its first Ext into $A^{N,0}$ zero,
so the family extends with the marked line fixed. The free PD
differential universal property extends the linear family uniquely
to $G_N\to A^K$. This proves the boundary and square filling
assertions, and hence the stated contractibility.
\end{proof}

\begin{theorem}[C17: \status{Proved under stated hypotheses}; finite-prime uniqueness]\label{sa08:C17}
Let $P$ be a finite set of rational primes and let
\[
 R=\Z[\ell^{-1}:\ell\notin P].
\]
For every commutative $R$-algebra $k$, the category of generalized
strict PD theories satisfying the uniform eventual conditions of
scans~117 and~119 has the following properties:
\begin{enumerate}
\item every pair has an admissible common source;
\item its localization at all morphisms is terminal.
\end{enumerate}
These conclusions remain valid if termwise $k$-flatness is required.
All comparison maps preserve both gradings, divided powers,
coefficients and augmentations. The theorem also holds in the
subcategory with $A^{w,q}=0$ for $q>w$, whenever that support
condition is included in the definition.
\end{theorem}
\begin{proof}
If $P$ is empty this is the rational case of
Statement~\ref{sa08:C09}. Assume $P$ is nonempty and put
$k_p=k\otimes_R\Z_{(p)}$, $k_{\Q}=k\otimes_R\Q$.

\paragraph{Arithmetic exactness.}
For every $R$-module $M$ there is a natural short exact sequence
\[
 0\longrightarrow M\longrightarrow
 M_{\Q}\oplus\prod_{p\in P}M_{(p)}
 \xrightarrow{\ \delta\ }(M_{\Q})^P\longrightarrow0,
 \qquad \delta(a,(b_p))=(b_p-a)_p.
\]
For $M=R$, the kernel is the intersection of the indicated
localizations inside $\Q$. Given a rational tuple $(y_p)$, choose
a common denominator $D$ and an integer $z$ satisfying
$z\equiv-Dy_p\pmod{p^{v_p(D)}}$ for each $p\in P$.
The Chinese remainder theorem supplies $z$; then $a=z/D$ and
$b_p=y_p+a$ give a preimage. All the coefficient modules are flat,
and the quotient $\Q^P$ is flat. The ring-level exact sequence is
therefore pure; tensoring it with $M$ proves the assertion without
any torsion restriction on $M$.

\paragraph{A natural arithmetic replacement.}
For a theory $A$ define
\[
 P(A)=\prod_{p\in P}A_{(p)},\qquad
 D(A)=P(A)_{\Q}=(A_{\Q})^P
\]
and, for $j=1,2$, form the strict pullback
\[
 J_j(A)=A_{\Q}\mathop{\times}_{D(A),e_0}
 D(A)^{I_j}\mathop{\times}_{e_1,D(A)}P(A).
\]
The natural map $\eta_A:A\to J_j(A)$ uses scalar extension and
constant paths. It preserves augmentations because $J_j(k)=k$:
constant coefficient modules acquire no extra simplicial paths,
and their arithmetic pullback is $k$.

Let $P_0D(A)$ denote paths whose initial endpoint is zero. Subtracting
the constant initial value defines an exact sequence in every
weight and cochain degree,
\[
 0\longrightarrow A\xrightarrow{\eta_A}J_j(A)
 \longrightarrow P_0D(A)\longrightarrow0.
\]
Surjectivity follows from arithmetic exactness: the final value of
any based path can be written $b-a$, with $a\in A_{\Q}$ and
$b\in P(A)$. Adding the constant path $a$ gives the required
lift. Based paths are simplicially contractible. For one edge use
the contraction of the interval to its initial vertex; for two
edges, subtract the initial value on each successive edge, which
identifies the based-path module with a direct sum of two
one-edge based-path modules. Consequently $\eta_A$ is a horizontal
weak equivalence in every weight and cochain degree.

Statement~\ref{sa08:C16}, applied componentwise, makes the based-path
kernel vertically acyclic in every fixed cochain range in all
large weights, uniformly in the simplex. Thus $J_j(A)$ is
admissible. If $A$ has flat terms, so does the based-path kernel
by the integral chain splitting in that statement, and hence so
does $J_j(A)$. Every map in this construction is functorial in $A$.

\paragraph{One glued admissible source.}
Given $A,B$, choose powers $N_p$ of $p$ large enough for
Statement~\ref{sa08:C09} to give
\[
 f_p:G_{N_p}\longrightarrow A_{(p)},\qquad
 h_p:G_{N_p}\longrightarrow B_{(p)}.
\]
Over $k_{\Q}$ choose a power $L$ of two, at least
$\max_{p\in P}N_p$ and sufficiently large, and write $C=G_L$.
The same lifting argument supplies maps
\[
 r_p:C\longrightarrow(G_{N_p})_{\Q},\qquad
 f_0:C\longrightarrow A_{\Q},\qquad
 h_0:C\longrightarrow B_{\Q}.
\]
Choose marked simplicial homotopies from $f_0$ to $(f_p)_{\Q}r_p$
and from $h_0$ to $(h_p)_{\Q}r_p$, using
Statement~\ref{sa08:C16}. Put
\[
 P_G=\prod_{p\in P}G_{N_p},\qquad D_G=(P_G)_{\Q},\qquad
 E=C\mathop{\times}_{D_G,e_0}D_G^{I_1}
       \mathop{\times}_{e_1,D_G}P_G.
\]
All rows of $C,P_G,D_G$ are acyclic in sufficiently large weights;
only finitely many thresholds occur for a fixed cochain range.
Endpoint restriction is then surjective. Its kernel and the
resulting pullback row are acyclic by the normalized Hom and
exact-sequence calculations of Statement~\ref{sa08:C16}. This
proves eventual horizontal acyclicity of $E$.

Vertically, in any fixed range and sufficiently large weights,
the relative path kernel has coefficient cohomology
$(S_{k_{\Q}}^w)^P$ in degree one only. The other two terms have
coefficient cohomology $S_{k_{\Q}}^w$ and
$\prod_pS_{k_p}^w$. The connecting map is the arithmetic difference
map $\delta$ above. Its kernel is $S_k^w$ and its cokernel is zero.
The long exact sequence therefore proves that $E$ resolves $S_k$
in the required eventual range, uniformly in the simplex.

The chosen homotopies give actual morphisms
\[
 E\longrightarrow J_2(A),\qquad E\longrightarrow J_2(B).
\]
For example, the first edge of the path in $J_2(A)$ is the chosen
homotopy from $f_0$ to $(f_p)_{\Q}r_p$, evaluated on the $C$
component. The second edge is the image under $f_p$ of the path
in $E$. Their values at the common vertex agree by the pullback
definition of $E$. This constructs a zigzag of admissible theories
\[
 A\longrightarrow J_2(A)\longleftarrow E
 \longrightarrow J_2(B)\longleftarrow B.
\]

We verify flatness of $E$ when it is required. The rational standard
row $G_N^{w,q}$ is acyclic unless $\lfloor w/N\rfloor=q$; in that
case it has homotopy $k_{\Q}$ in degree $q$. This follows from
$G_N$'s finite-free description: a positive divided-power factor
of $\Phi$ is contractible, while
$N(\Lambda^q\Psi)=k_{\Q}[q]$, directly from the exterior basis
of $\Psi_n=k_{\Q}^n$. Over $\Q$ the normalized complexes split;
base change shows that their terms and cycles are free over
$k_{\Q}$. The kernel $K$ of endpoint restriction on $D_G^{I_1}$
is therefore flat. Indeed the integral relative prism complex is
$\Z[1]$ plus identity disks, so its Hom uses $Z_1(ND_G)$ and
normalized terms, all flat.

The cokernel of endpoint restriction is the constant module
$\pi_0(D_G)$. It is zero for $q>0$. For $q=0$ its $p$-component
is $k_{\Q}$ exactly when $w<N_p$. In these weights $C^{w,0}=k_{\Q}$
because $w<L$. Thus the kernel $V$ of
$C\oplus P_G\to\pi_0(D_G)$ is the arithmetic kernel for the subset
\[
 P_w=\{p\in P:w<N_p\},
\]
plus the remaining flat terms of $P_G$. That arithmetic kernel
is $k$ localized by the primes in $P\setminus P_w$, hence is flat.
If $P_w$ is empty, $V=C\oplus P_G$ is already flat. The exact
sequence $0\to K\to E\to V\to0$ now proves flatness in all weights
and cochain degrees, including those outside the eventual range.

\paragraph{An actual common source.}
We turn the displayed zigzag into a common source without asserting
strict equalization. For $\eta_A:A\to J_2(A)$ take its simplicial
mapping-path factorization
\[
 \widetilde A=A\mathop{\times}_{J_2(A),e_0}J_2(A)^{I_1}
 \xrightarrow{e_1}J_2(A).
\]
Projection $\widetilde A\to A$ has a constant-path section and
based-path kernel. Thus it is a horizontal weak equivalence in
every row. Since $\eta_A$ is also a horizontal weak equivalence,
$e_1$ is a trivial fibration of simplicial modules in every row:
endpoint restriction is a simplicial fibration, and the normalized
mapping-fiber sequence computes its homotopy groups. In particular
$e_1$ is degreewise surjective with an acyclic kernel. This can
also be checked directly: the obstruction to a path between two
simplices is their difference in $\pi_0$, and $\eta_A$ induces
an isomorphism on every homotopy group.

The relative comparison in Statement~\ref{sa08:C16} makes
$\widetilde A$ admissible. Its kernel over $J_2(A)$ is vertically
acyclic in every fixed range in all sufficiently large weights,
since the map between the two resolutions preserves $S_k$.
Construct $\widetilde B$ similarly and set
\[
 H=E\mathop{\times}_{J_2(A)\times J_2(B)}
       (\widetilde A\times\widetilde B).
\]
The projection $H\to E$ is degreewise surjective. Its kernel is
the direct sum of the two preceding kernels, hence is horizontally
acyclic in every row and vertically eventually acyclic. The exact
sequences prove that $H$ is admissible. Its projections give the
required maps $H\to A$ and $H\to B$.

If flatness is imposed, $\widetilde A$ is degreewise the direct
sum of $A$ and based paths in $J_2(A)$, so is flat. The kernel of
its surjection onto the flat $J_2(A)$ is flat as well. The same
holds for $B$, and the extension defining $H$ preserves flatness.
This proves the common-source assertion in both categories.

\paragraph{Parallel maps and localization.}
Let $u,v:A\rightrightarrows B$. Use the same construction with
$N_p,L$ also beyond the degree-zero acyclicity thresholds of $B$.
Write $H_p$ for the homotopy from $f_0$ to $(f_p)_{\Q}r_p$ that
defines $h:E\to J_2(A)$. Choose local marked homotopies $K_p$
from $u_{(p)}f_p$ to $v_{(p)}f_p$, and a rational marked homotopy
$K_0$ from $u_{\Q}f_0$ to $v_{\Q}f_0$.

For each $p$, the four sides
\[
 u_{\Q}H_p,\qquad v_{\Q}H_p,\qquad
 K_0,\qquad (K_p)_{\Q}r_p
\]
form a compatible square of marked maps from $C$ into $B_{\Q}$.
Statement~\ref{sa08:C16} fills this square. The filled square
supplies the first edge of a homotopy
$E\to J_2(B)^{I_1}$; applying $K_p$ to the path in $E$ supplies
its second edge. The two edges agree at their common vertex.
Its endpoints are exactly $J_2(u)h$ and $J_2(v)h$.

The target $J_2(B)^{I_1}$ is admissible. Its two evaluations become
equal after localization, since their common constant-path section
is inverted. The map $h$ and the natural maps $\eta_A,\eta_B$ are
also inverted. Naturality of $\eta$ consequently gives $u=v$ in
the localized category.

Finally, common sources for pairs give a common source for every
finite list of objects, by iteration. For a finite zigzag loop,
choose such a source mapping to every object on the loop. For
each arrow, its two composites from that source are parallel
actual morphisms, hence coincide after localization by the previous
paragraph. The loop therefore telescopes to the identity. The
category is connected by the constructed common sources, so its
localized groupoid is terminal.
\end{proof}

The remaining coefficient-ring gate has infinitely many noninvertible
rational primes, already $k=\Z$. The finite-prime proof uses finite
maxima of local thresholds. Statement~\ref{sa08:C12} rules out applying
it unchanged to an unbounded family of the local sources $G_{N_p}$.
